% LaTeX source for book ``代数学方法'' in Chinese
% Copyright 2018  李文威 (Wen-Wei Li).
% Permission is granted to copy, distribute and/or modify this
% document under the terms of the Creative Commons
% Attribution 4.0 International (CC BY 4.0)
% http://creativecommons.org/licenses/by/4.0/

% 《基础代数几何 I：射影空间中的簇》中译本 — 导言
% 底本: I.R. Shafarevich, Basic Algebraic Geometry 1, 3rd ed., Springer, 2013.
	\chapter*{导言}

\section*{概观}

本册是沙法列维奇 (\trShafarevich)《基础代数几何》第一卷《射影空间中的簇》的中文译本\cite{shafarevich2013}. 原书以经典代数几何 (概形理论之前) 的语言展开: 从平面代数曲线与仿射空间中的闭子集讲起, 依次处理局部性质, 除子与微分形式, 相交数, 最后以代数附录总结所需的交换代数工具. 与同题材的现代教材相比, 原书以例子与几何直观见长, 且习题丰富.

\section*{简要说明}
\paragraph*{旨趣}
本译本的底本是 Miles Reid 的英译本 (Springer, 第三版, 2013). 翻译约定如下:
\begin{enumerate}
	\item 数学正确性优先: 正文记号悉依原文; 原书疑有笔误处照录原文, 并以``译者注''说明, 不擅自改写数学内容.
	\item 术语与人名译法参照《The Rising Sea》中译本术语表, 全书统一.
	\item 公式与定理编号按章编排, 与原书一致, 便于对照.
	\item 插图系据原图以 TikZ 重绘, 图号与原书对应.
	\item 本导言及各章内容提要为译本所增补, 原书无此内容.
\end{enumerate}

\section*{内容提要}
以下简介各章的内容.
\begin{asparadesc}
	\item [第一章\ 基本概念] 从平面代数曲线的例子入手, 引入闭子集, 正则函数, 有理函数, 拟投影簇, 有限映射与维数等基本观念.

	\item [第二章\ 局部性质] 奇异点与切空间, 幂级数展开, 双有理映射的结构 (胀开), 正规簇与映射的奇异性.

	\item [第三章\ 除子与微分形式] 除子理论, 曲线上的除子, 平面三次曲线与群律, 代数群, 微分形式, 曲线上的 Riemann--Roch 定理.

	\item [第四章\ 相交数] 相交数的定义与基本性质, 其应用 (含三次曲面), 曲面上的双有理映射, 奇点.

	\item [代数附录] 线性与双线性代数, 多项式, 拟线性映射, 不变量, 域, 交换环, 唯一分解, 整元素与模的长度.
\end{asparadesc}

\paragraph*{附记}


\section*{致谢}
本译本采用李文威老师的 AJbook 文档类 (\LaTeX{} 模板), 感谢李文威老师在技术上的无私分享.



\let\cleardoublepage\clearpage
% % % % % % % % % %
\part{射影空间中的簇}

\let\cleardoublepage\clearpage
